%% Chapter 19 ----------------------------------------------------------------
\chapter{Coronal Magnetic Field Modeling (PFSS)}
\label{ch:pfss}
\index{PFSS}\index{potential field source surface}

An EUV image shows where the coronal plasma is, not where the magnetic field\index{magnetic field} is. The
\gui{Potential Field (PFSS)} card of Solar Image Analysis extrapolates the coronal magnetic
field from a photospheric magnetogram with the potential-field source-surface model
\parencite{schatten1969,altschuler1969} and draws the resulting field lines over the disk
(Figure~\ref{fig:pfss}), so that the modeled loop topology and coronal holes can be compared
with the image. The card is shown for disk imagers, magnetographs and synoptic maps.

\screenshot{solar/pfss_overlay}{PFSS field lines over an AIA image.}{fig:pfss}{An AIA 193 Å image with PFSS open field lines (red/blue), gray closed loops and open-field boundaries drawn, and the Potential Field (PFSS) card visible with its settings and status line.}

\begin{background}[Background: the PFSS model]
  The model assumes a current-free (potential) field between the photosphere and a spherical
  \emph{source surface} a few solar radii out, where the field is forced to be radial to
  mimic the solar wind dragging it open. Field lines that reach the source surface are
  \emph{open} and feed the solar wind; the others are \emph{closed} loops. The open-field
  footpoint regions are the model's coronal holes, and the polarity-inversion line on the
  source surface marks where the heliospheric current sheet leaves the model.
\end{background}

\begin{caution}
  AIA images carry no magnetic-field information, so the boundary condition comes from a
  \emph{synoptic} magnetogram --- a full-Sun map assembled from central-meridian observations
  over a whole solar rotation. The model therefore describes the global field around your
  observation, never the instantaneous field, and the far side is always at least two weeks
  old. The card shows how far the map's date is from the displayed frame; that offset is the
  number to judge the result by.
\end{caution}

\section{Settings}

\begin{table}[!htbp]
  \centering
  \caption{Settings of the \gui{Potential Field (PFSS)} card.}
  \label{tab:pfss}
  \small
  \begin{tabularx}{\linewidth}{@{}P{0.19\linewidth}L@{}}
    \toprule
    \thead{Setting} & \thead{Description} \\
    \midrule
    Magnetogram & The boundary condition (Section~\ref{sec:pfss-magnetograms}). For \gui{GONG ADAPT}\index{PFSS!ADAPT}, \gui{Realization} (0--11) chooses one of the twelve flux-transport realizations; for \gui{Local FITS file...}, \gui{Browse…} selects the file. The line below names the loaded map and its time offset from the frame. \\
    Source surface\index{PFSS!source surface} & Height at which the field is forced radial (1.5--10~\(R_\odot\); default 2.5~\(R_\odot\)). Lowering it opens more field. \\
    Radial cells\index{PFSS!radial cells} & Radial grid resolution between the photosphere and the source surface (10--150; default 35). More is smoother and slower. \\
    Seeds & Where field lines start (Section~\ref{sec:pfss-seeds}). \\
    Density & Seed density (4--64; default 12, about 150 field lines). It controls both speed and legibility; the hint beside it estimates the number of lines. \\
    \bottomrule
  \end{tabularx}
\end{table}

\subsection{Magnetograms}
\label{sec:pfss-magnetograms}
\index{magnetogram!synoptic}

\begin{description}
  \item[GONG synoptic] the ground-based GONG synoptic\index{PFSS!GONG} map \parencite{harvey1996}, updated daily
    and available for any date from 2006 on. No registration is needed; this is the usual
    choice.
  \item[HMI synoptic (JSOC)] SDO's own synoptic map, instrumentally consistent with AIA.
    Searching is free, but downloading stages a JSOC\index{JSOC} export and needs the registered e-mail
    set in the \gui{Data Source} card. The series is published about one Carrington rotation
    behind, so the most recent weeks are not covered.
  \item[GONG ADAPT] maps from the ADAPT flux-transport model, which estimates the unobserved
    far side instead of leaving it weeks old and usually improves open-field predictions. One
    file holds twelve realizations, which differ mostly on the far side.
  \item[Local FITS file...] a synoptic map you already have. It must be a full-Sun Carrington
    map; plate-carrée maps are reprojected to equal-area automatically.
\end{description}
Downloaded magnetograms are cached, so changing a model parameter and solving again does not
download the map again. Unmeasured pixels, such as the polar gaps in GONG maps, are filled
with zero so that the solution can be computed; their number is reported.

\subsection{Seeds}
\label{sec:pfss-seeds}

\begin{description}
  \item[Uniform grid] an equal-area sample over the visible hemisphere.
  \item[Open-field regions] a coarse survey is traced first, and then lines are started only
    where the field is open. This shows coronal-hole connectivity without closed loops
    dominating the picture.
  \item[Current crop] concentrates the lines inside the crop rectangle.
  \item[Clicked points] traces one line from each point clicked on the disk. This needs the
    PyQtGraph\index{PyQtGraph} renderer.
\end{description}
Tracing costs about 1.6~ms per field line, so a solution takes a few seconds at most, but
beyond about 150 lines the field lines start to hide the image they are compared with. Prefer
\gui{Open-field regions}\index{PFSS!seeds} to a high density when coronal-hole connectivity is what you want.

\section{Computing and displaying the model}

\gui{Compute PFSS}\index{PFSS!compute} downloads the magnetogram if needed, solves the potential field and traces
the field lines in the background; \gui{Clear} removes the overlay and the solution. The
\gui{Show} options toggle the layers without recomputing:
\begin{description}
  \item[Open field] lines reaching the source surface, colored by the sign of the radial field
    at their footpoint: red outward, blue inward.
  \item[Closed loops] lines returning to the photosphere, drawn in gray.
  \item[Open-field boundaries] outlines of the open-field footpoint regions, for comparison
    with the dark coronal holes in 193 or 211~Å images.
  \item[Near side only] (on by default) hides lines rooted on the far hemisphere. Such lines
    are not hidden behind the disk once they rise above the limb, so they would otherwise be
    drawn as arcs outside the disk with no visible footpoint. Clear it for the geometrically
    complete picture.
\end{description}
Stepping through a sequence re-projects the same solution onto each frame instead of solving
again. This is correct, because the synoptic map is valid across the rotation while the
observer's view changes from frame to frame.

\section{Diagnostics}

\gui{Diagnostics…}\index{PFSS!diagnostics} opens a four-panel window that documents the solution
(Figure~\ref{fig:pfss-diagnostics}):
\begin{enumerate}[label=\arabic*.]
  \item the input magnetogram (\(B_r\)) with its provenance;
  \item the radial field at the source surface with its neutral line, where the heliospheric
    current sheet leaves the model --- the model's most directly testable prediction;
  \item the open (red and blue) and closed (dark) footpoint map in Carrington projection;
  \item the solution statistics, including the open flux, the open-area fraction and the
    number of filled pixels.
\end{enumerate}
\gui{Save Figure…} saves all four panels as one image (PNG, PDF or SVG).

\screenshot{solar/pfss_diagnostics}{The PFSS diagnostics window.}{fig:pfss-diagnostics}{The PFSS diagnostics window showing the four panels: input magnetogram, source-surface Br with neutral line, open/closed footpoint map, and solution statistics.}

\begin{note}
  PFSS modeling uses the \texttt{sunkit-magex}\index{sunkit-magex} package \parencite{stansby2020}, which is
  included in the installers. When the application is run from source without it, the card
  remains visible but disabled and shows the command to install it. Sessions store the PFSS
  settings, not the solution (Section~\ref{sec:solar-sessions}).
\end{note}
